% English translation, made from our own transcription (original.tex), not from any existing
% translation. Every math expression, symbol, \tag, \origpage{N} marker and \emph run is reproduced
% unchanged; only the prose is translated.
%
% PROSE INSIDE FORMULA TRANSLATED (HOUSESTYLE R21):
% The connective in the display on p. 1148 (\int \frac{B\,dx - A\,dy}{f'_{z}} \quad\text{et}\quad ...)
% is translated from \text{et} to \text{and}.
%
% QUOTATION MARKS ARE NOT TRANSLATED (HOUSESTYLE R18, R22):
% The paragraph-initial « / » of the Comptes rendus and the terminal » at the end of p. 1149 are
% typographic features of the source edition and stand exactly where the transcription puts them,
% set tight against the text.
%
% TERMINOLOGY:
%   * intégrales de différentielles totales = "integrals of total differentials"
%   * première espèce = "first kind"
%   * fonctions uniformes = "single-valued functions"
%   * fonctions quadruplement périodiques = "quadruply periodic functions"
%   * surface d'ordre m = "surface of order m"
%   * singularités ordinaires = "ordinary singularities"
%   * point quelconque = "arbitrary point"
%   * multiples des périodes = "multiples of the periods"
%   * conditions nécessaires et suffisantes = "necessary and sufficient conditions"
%   * courbe double = "double curve"
%   * déterminant ... n'est pas identiquement nul = "determinant ... is not identically zero"
%   * points doubles isolés = "isolated double points"
%   * genre = "genus"
%   * paires de périodes simultanées = "pairs of simultaneous periods"
%   * équations aux différentielles totales = "total differential equations"
%   * section plane quelconque = "arbitrary plane section"
%   * intégrales abéliennes = "abelian integrals"
\documentclass{article}
\usepackage{readmasters}
\begin{document}

\origpage{1147}
\section*{Mathematical Analysis. --- On integrals of total differentials and on a class of algebraic surfaces. Note by M.~E.~Picard, presented by M.~Hermite.}

«In a recent Communication (\emph{Comptes rendus}, meeting of 1~December 1884), I stated a fundamental proposition concerning integrals of algebraic total differentials that always remain finite, integrals which I designate for this reason by the name of \emph{integrals of the first kind}. The applications of this theory, which I am currently developing in an extensive Memoir, may perhaps present some
\origpage{1148}interest; among these applications, one of the simplest concerns algebraic surfaces whose coordinates can be expressed by single-valued quadruply periodic functions of two parameters. I request permission to state here a general proposition concerning these surfaces.

»Given a surface of order $m$
\[
f(x, y, z) = 0,
\]
possessing only the ordinary singularities considered in the cited Note, let us investigate how one can recognize whether the coordinates $x$, $y$, $z$ of an arbitrary point can be expressed by quadruply periodic functions of two parameters, and this in such a manner that to an arbitrary point of the surface there corresponds only a single system of values of the two parameters, disregarding multiples of the periods. The \emph{necessary and sufficient} conditions for this to be the case can be formulated in the following manner:

»\emph{The proposed surface has a double curve of order $\frac{m(m - 4)}{2}$ and it possesses two integrals of total differentials of the first kind}
\[
\int \frac{B\,dx - A\,dy}{f'_{z}} \quad\text{and}\quad \int \frac{B_{1}\,dx - A_{1}\,dy}{f'_{z}},
\]
\emph{for which the determinant $AB_{1} - A_{1}B$ is not identically zero.}

»The number $m$ is necessarily even, and it follows from the preceding conditions that the surface has no \emph{isolated} double points; furthermore, the \emph{genus} of the surface is equal to unity.

»The two preceding integrals will have \emph{four} pairs of simultaneous periods, and the two total differential equations
\[
\begin{aligned}
\frac{B\,dx - A\,dy}{f'_{z}} &= du, \\
\frac{B_{1}\,dx - A_{1}\,dy}{f'_{z}} &= dv
\end{aligned}
\]
will give for $x$ and $y$ single-valued, quadruply periodic functions of $u$ and $v$.

»We have shown how one can recognize whether a surface admits integrals of the first kind, and find them when they exist; the proposed problem is therefore completely solved.

»An arbitrary plane section of the surface is of degree $m$ and of
\origpage{1149}genus $\frac{m + 2}{2}$, and this curve presents an interesting peculiarity from the point of view of the reduction of the number of periods of the abelian integrals corresponding to it.

»It follows, indeed, immediately from the preceding results that, among these abelian integrals, \emph{there are two that have only four pairs of corresponding periods}.»

\end{document}
